% Generated draft: requires mathematical and source review.
\documentclass[11pt]{article}
\usepackage[margin=25mm]{geometry}
\usepackage{fontspec}
\setmainfont{Latin Modern Roman}
\usepackage{amsmath,amssymb,mathtools,unicode-math}
\setmathfont{Latin Modern Math}
\usepackage{enumitem,needspace,xurl,hyperref}
\setlength{\parindent}{0pt}
\setlength{\parskip}{5pt}
\setlength{\emergencystretch}{3em}
\newcommand{\N}{\mathbb{N}}
\newcommand{\Z}{\mathbb{Z}}
\newcommand{\Q}{\mathbb{Q}}
\newcommand{\R}{\mathbb{R}}
\newcommand{\C}{\mathbb{C}}
\newcommand{\F}{\mathbb{F}}
\newcommand{\A}{\mathbb{A}}
\newcommand{\PP}{\mathbb{P}}
\newcommand{\E}{\mathbb{E}}
\newcommand{\eps}{\varepsilon}
\newcommand{\dd}{\,\mathrm{d}}
\begin{document}
\section*{UnsolvedMath ID: 700712}
\textbf{Source problem code:} catalogue entry
\subsection{Definitions and mathematical statement}
Let \(F\) be an absolutely free group, meaning a free group in the variety of all groups, and let \(X\) be a fixed free generating set. Use the commutator convention \([a,b]=a^{-1}b^{-1}ab\). Define the lower central series by
\[
\gamma_1(F)=F,\qquad \gamma_{k+1}(F)=[\gamma_k(F),F]\quad(k\geq 1),
\]
where \([H,K]\) denotes the subgroup generated by all \([h,k]\) with \(h\in H\) and \(k\in K\).

Interpret basic commutators using a fixed Hall ordering: the elements of \(X\) are basic commutators of weight \(1\); higher weights follow lower weights in the ordering; and \([u,v]\) is a basic commutator of weight \(\operatorname{wt}(u)+\operatorname{wt}(v)\) when \(u,v\) are basic commutators, \(u>v\), and, if \(u=[u_1,u_2]\), then \(u_2\leq v\). Fix an order among the basic commutators of each weight. Write \(\mathcal B_n(X)\) for the resulting set of basic commutators of weight exactly \(n\), evaluated in \(F\). For \(S\subseteq F\), its normal closure in \(F\) is
\[
\langle\!\langle S\rangle\!\rangle_F
=\bigl\langle g^{-1}sg : g\in F,\ s\in S\bigr\rangle.
\]

\textbf{Question (attributed in the supplied statement to C. Sims).} For every positive integer \(n\), is
\[
\gamma_n(F)=\langle\!\langle\mathcal B_n(X)\rangle\!\rangle_F?
\]
The question is posed for a free group with a fixed free basis. The supplied source imposes no explicit restriction on the rank of \(F\), and does not specify the Hall ordering; the definitions above make the standard Hall interpretation explicit. The precise intended scope over ranks and choices of basic commutators requires source review.

The supplied record reports that D. Jackson announced a positive answer for \(n\leq 5\), here understood as \(1\leq n\leq 5\). This is retained as a reported announcement, without independent verification or an assertion that a proof was published.
\subsection{Short English statement}
For a free group with fixed free generators, does the normal closure of the basic commutators of weight exactly \(n\) equal the \(n\)th term of its lower central series, for every positive integer \(n\)? The source reports a positive answer announced by D. Jackson for \(1\leq n\leq 5\); the intended rank scope and commutator conventions need review.
\subsection{Sources}
\begin{itemize}
\item Local UnsolvedMath registry: permanent ID \(700712\), exact source code \texttt{catalogue entry}, title \emph{Kourovka Notebook Problem 11.19}. The supplied record labels the problem open; this label is supplied registry data, not an independently verified literature assessment.
\item \emph{The Kourovka Notebook (Unsolved Problems in Group Theory)}, original problem 11.19, 11th issue (1990); supplied edition reference: 21st issue edition (2026), \texttt{arXiv:1401.0300v48}. The supplied collection description names E. I. Khukhro and V. D. Mazurov as editors.
\item Attribution: the record lists A. Gaglione and D. Spellman as proposers, with proposed year 1990, while the statement itself begins with an attribution to C. Sims. Their respective roles are not explained in the supplied record.
\item The supplied statement reports D. Jackson's announcement of a positive answer for \(n\leq 5\). No separate bibliographic reference or proof citation for that announcement is supplied.
\end{itemize}
\end{document}
